Neural algorithmic reasoning asks whether networks can learn to execute classical algorithms and keep working on inputs larger than those seen in training \cite{velikovi2021neural}. Graph algorithms are the standard testbed: the CLRS benchmark trains on small instances and tests on larger ones \cite{velikovi2022clrs}, and SALSA-CLRS pushes this to large sparse graphs \cite{minder2023salsa}. Two design choices are repeatedly credited for the success of message-passing networks on such tasks. The first is \emph{algorithmic alignment}: the architecture should mirror the structure of the algorithm, which lowers sample complexity \cite{xu2019what} and, for shortest paths, makes max (or min) aggregation a natural match for the relaxation step \cite{xu2020how,dudzik2022graph}. The second is \emph{step-wise supervision}: training the network to reproduce the algorithm's intermediate states, as in neural execution of graph algorithms \cite{velikovi2019neural} and the hints of CLRS.

These two claims are usually evaluated inside larger systems that change several things at once. Here we isolate them for one algorithm, Bellman-Ford, in a $2\times2$ design (sum vs.\ max aggregation, with vs.\ without step supervision) plus an MLP that sees the flattened adjacency matrix. Networks are trained on graphs of at most 8 nodes and tested on 8, 16, 32 and 64 nodes, on two graph families that differ in whether node degree stays constant or grows with $n$; degree shift is a known cause of size-generalisation failure in GNNs \cite{yehudai2020from}. We state four falsifiable hypotheses (Section~\ref{sec:setup}) before running and report all seeds.

\textbf{Contributions.} (i) A controlled, fully reproducible CPU-scale factorial study of aggregation and step supervision for size extrapolation of Bellman-Ford; (ii) a negative finding: step supervision did not help under our fixed budget, and sum aggregation with steps produced non-finite or astronomically large errors on growing-degree graphs at 64 nodes; (iii) a seed-level analysis showing that mean errors at 64 nodes are dominated by a minority of failing seeds, and a test-time step sweep of the two step-supervised systems (three seeds each) in which extra message-passing steps raise the error instead of leaving it unchanged. We claim nothing beyond this scale: one algorithm, small networks (the MPNN has 12,770 parameters and the MLP 557,376), five seeds, no hyperparameter search.
